% =============================================================================
% conclusiones.tex — Conclusiones y Recomendaciones
%
% Norma §1.2.3: Las conclusiones deben redactarse en relación directa con
% los objetivos específicos, en forma breve, clara, objetiva y concisa.
%
% Norma §1.2.4: Las recomendaciones se formulan en base a las conclusiones,
% considerando aspectos susceptibles de mejora en trabajos futuros.
%
% Ambas secciones NO son capítulos (no llevan numeración de capítulo).
% =============================================================================

% ── Hoja separadora: CONCLUSIONES ────────────────────────────────────────────
\seccioncuerpo{CONCLUSIONES}

% Redactar una conclusión por cada objetivo específico, usando conectores
% ordinales: PRIMERA, SEGUNDA, TERCERA, CUARTA, QUINTA, SEXTA.

\textbf{PRIMERA:}

\textbf{SEGUNDA:}

\textbf{TERCERA:}

\textbf{CUARTA:}

\textbf{QUINTA:}

\textbf{SEXTA:}

% ── Hoja separadora: RECOMENDACIONES ─────────────────────────────────────────
\seccioncuerpo{RECOMENDACIONES}

% Redactar recomendaciones derivadas de las conclusiones.
% Deben considerar aspectos susceptibles de mejora e implicaciones en trabajos
% futuros. Ordenarlas por prioridad.
